% Source: https://artofproblemsolving.com/community/c3h2505356_harmonic_mean_property
% Stable identity: aops-topic-2505356-q01
\textbf{(a)} We compute
\[h-a=\dfrac{2ab-a^2-ab}{a+b}=\dfrac{a(b-a)}{a+b},\qquad b-h=\dfrac{ab+b^2-2ab}{a+b}=\dfrac{b(b-a)}{a+b},\]
so $\dfrac{h-a}{b-h}=\dfrac ab$. Since $\dfrac ab>0$, the differences $h-a$ and $b-h$ have the same sign, so $h$ lies between $a$ and $b$, and it splits the interval from $a$ to $b$ in the ratio $a:b$; that is, $h=a+\dfrac{a}{a+b}(b-a)$. For $12$ and $24$: $h=12+\dfrac{12}{36}\cdot12=\boxed{16}$ (check: $\dfrac{2\cdot12\cdot24}{36}=16$).

\textbf{(b)} The times are $t_a=\dfrac da$ and $t_b=\dfrac db$, so the average speed is
\[\dfrac{2d}{\dfrac da+\dfrac db}=\dfrac{2ab}{a+b}=h.\]
An average speed is a \emph{time}-weighted average of the speeds: $h=\dfrac{a\,t_a+b\,t_b}{t_a+t_b}$. A weighted average divides the interval between $a$ and $b$ inversely to the weights: $\dfrac{h-a}{b-h}=\dfrac{t_b}{t_a}$. Equal distances mean $t_a:t_b=\dfrac1a:\dfrac1b=b:a$, so $\dfrac{h-a}{b-h}=\dfrac ab$, which is the property in (a). The runner spends more time at the slower speed, so the average is pulled towards it. For example, $12$ miles at $4$ miles per hour and $12$ miles at $6$ miles per hour take $3$ and $2$ hours, and the average speed is $4+\dfrac{2}{5}\cdot2=4.8$ miles per hour.
